\chapter*{不定积分每日一题}

\refer{《积分的方法与技巧》金玉明\ 顾新身\ 毛瑞庭\ /\ 编著}

\exercise{求积分$\int{\frac{\d x}{x+\sqrt{x^2-5x+6}}}$}%在花括号内输入题目

\teacher{陆丽芬}%在花括号内输入你的名字

\postdate{2022年12月18日}%在花括号内输入你想要发布那一天的日期

\bigresult{
    \formula{
        \begin{aligned}
            \int{\frac{2t\d t}{\qty(t+1)^2\qty(t-1)\qty(2t-3)}}=&\frac{1}{10}\cdot\frac{1}{2\sqrt{x^2-5x+6}+2x-5}+\frac{24}{50}\ln{\qty(\sqrt{x^2-5x+6}+x)}\\
            &+\frac{1}{50}\ln{\qty(5-2x-2\sqrt{x^2-5x+6})}+C
        \end{aligned}
    }
}%在花括号内输入本题结果

\analysis{
    本题可以使用欧拉换元积分法，对于欧拉换元积分法详细内容可以看昨日的每 日一题(12.17).下面直接操作一下.

    观察结构，我们发现
    \formula{x^2-5x+6=\qty(x-2)\qty(x-3)}

    因此，这属于第三种替换情况，可令
    \formula{\sqrt{x^2-5x+6}=t\qty(x-2)}

    两边同时平方可得
    \formula{x^2-5x+6=t^2\qty(x-2)^2}

    而同时，我们有
    \formula{x^2-5x+6=\qty(x-2)\qty(x-3)}

    联立后化简可得
    \formula{x-3=t^2\qty(x-2)}

    整理一下可得
    \formula{t^2=\frac{x-3}{x-2}=1-\frac{1}{x-2}}
    \formula{x=\frac{1}{1-t^2}+2=\frac{2t^2-3}{t^2-1}}

    对上述式子两边同时求导化简可计算出$\d x$，操作如下
    \formula{t^2=1-\frac{1}{x-2}}
    \formula{2t\d t=\frac{1}{\qty(x-2)^2}\d x}
    \formula{\d x=2t\qty(x-2)^2=\frac{2t\d t}{\qty(1-t^2)^2}}

    有了这些我们就可以去直接代入计算原积分了
    \formula{
        \begin{aligned}
            \int{\frac{\d x}{x+\sqrt{x^2-5x+6}}}&=\int{\frac{1}{\dfrac{2t^2-3}{t^2-1}+t\qty(x-2)}\cdot\frac{2t}{\qty(1-t^2)^2}\d t}\\
            &=\int{\frac{2t\d t}{\qty(2t^2-3)\qty(t^2-1)+t\cdot\frac{1}{1-t^2}\cdot\qty(t^2-1)^2}}\\
            &=\int{\frac{2t\d t}{\qty(2t^2-3)\qty(t^2-1)+t\qty(1-t^2)}}\\
            &=\int{\frac{2t\d t}{\qty(2t^2-t3)\qty(t^2-1)}}\\
            &=\int{\frac{2t\d t}{\qty(t+1)^2\qty(t-1)\qty(2t-3)}}
        \end{aligned}
    }
    到这，我们就已经把原本复杂的根式积分化为了有理分式积分，下面可以使用待定系数法分解被积函数.
    \formula{\frac{2t}{\qty(t+1)^2\qty(t-1)\qty(2t-3)}=\frac{At+B}{\qty(t+1)^2}+\frac{C}{\qty(t+1)^2}+\frac{D}{t-1}+\frac{E}{2t-3}}
    通分下可得
    \formula{
        \addtolength{\jot}{-1em}
        \begin{aligned}
            2t=&\qty(At+B)\qty(t-1)\qty(2t-3)+C\qty(t+1)\qty(t-1)\qty(2t-3)\\
            &+D\qty(t+1)^2\qty(2t-3)+E\qty(t+1)^2\qty(t-1)
        \end{aligned}
    }\addtolength{\jot}{-1em}
    我们可以通过特殊值代入、两边系数对比等办法求解这些系数

    比如代入$t=1$，可得$D=\frac{1}{2}$；代入$t=-1$，可得$A-B=\frac{1}{5}$；代入$t=\frac{3}{2}$，可得$E=\frac{24}{25}$；代入$t=0$，可得$B+C=-\frac{9}{50}\cdots\cdots$

    最后可得，不要忘记代入哦$\sim$
    \formula{
        \begin{aligned}
            \int{\frac{2t\d t}{\qty(t+1)^2\qty(t-1)\qty(2t-3)}}=&\int{\qty[-\frac{1}{5\qty(t+1)^2}+\frac{1}{50\qty(t+1)}-\frac{1}{2\qty(t-1)}+\frac{24}{25\qty(2t-3)}]\d t}\\
            =&\frac{1}{50}\qty[\frac{10}{t+1}+\ln{\qty(t+1)}-25\ln{\qty(1-t)}=24\ln{\qty(3-2t)}]+C\\
            =&\frac{1}{5\sqrt{\dfrac{x-3}{x-2}}+5}+\frac{1}{50}\ln{\qty(\sqrt{\frac{x-3}{x+2}}+1)}-\frac{1}{2}\ln{\qty(1-\sqrt{\frac{x-3}{x-2}})}\\
            &+\frac{24}{50}\ln{\qty(3-2\sqrt{\frac{x-3}{x-2}})}+C
        \end{aligned}
    }
    最后的化简就不展示了，化简不要求，到这答案也足够了.
}%在花括号内输入分析